% !TeX root = ../main.tex
\section{Discussion}
\label{sec:discussion}

\subsection{Insights on Times-Intersect and its Measurement}

In Fig.~\ref{fig:tintDistrib}, across all beam sequences, $t_{\textnormal{int}}$ decreased nonlinearly with stepped increases shear modulus. $t_{\textnormal{int}}$ from materials with shear moduli under 15 kPa were distributed more broadly, and adjacent boxes were easily distinguishable (Hedge's $g > 0.5$). On the other hand, distributions of $t_{\textnormal{int}}$ became tighter in stiffer materials, and differences between adjacent elasticites exhibited smaller effect sizes. Therefore, $t_{\textnormal{int}}$ is more sensitive to relative changes in softer tissues than in stiffer ones.

Figure~\ref{fig:tintStat} quantifies these differences, showing that effect sizes and statistical separability were largest for soft materials and diminished for stiffer ones. Beam configuration also played a role: wider push beams, larger differences in F-number (e.g. F/1.5 v. F/5.0), and simultaneous rather than sequential tracking improved separability. However, these advantages decreased as stiffness increased, with statistical significance sometimes lost in the upper modulus ranges.

Together, these findings suggest that $t_{\textnormal{int}}$ encodes stiffness in a nonlinear manner that favors the discrimination of softer tissues. The smaller effect sizes for pairs of $t_{\textnormal{int}}$ from stiffer materials highlights a fundamental limitation in resolving stiffer materials. Since scatterers shear less and recover from ARF excitations more quickly in stiffer materials: (i) times-intersect are less distinguishable since the paired displacement profiles, themselves, are more similar; (ii) the ability to detect times-intersect is limited by the pulse repetition interval (PRI) of track pulses, and (iii) smaller displacement magnitudes become more difficult to to detect given displacement estimation variance in stiffer materials~\cite{mcaleavey_Estimates_2003,dhanaliwala_Assessing_2012}.

% OMITTED FROM SPRING 2025 DRAFT:
%
% Fig.~\ref{fig:tintStat} shows that, in all evaluated beam sequences, times-intersect are distributed in a nonlinear manner as a function of the true elasticity of the ARF-interrogated \emph{in silico} medium. Times-intersect are distributed such that lower shear modulus values are associated with higher values, whereas stiffer materials are associated with lower values. The nonlinear relationship specifically shows that the differences in times-intersect are more pronounced for softer materials while times-intersect distributions from stiffer materials are closer together. This pattern appears consistent with the pattern from ARFI imaging~\cite{nightingale_Acoustic_2002} where softer materials displace and recover from ARF excitations more slowly than stiffer materials. In ARFI, the rate of scatterer shearing is faster in stiffer materials to the point where the ability to observe ARF-induced displacements is limited by the pulse repetition interval (PRI), the inverse of the PRF; displacement responses that occur and recover more quickly cannot be detected ultrasonically if the PRF is too low. This is consistent with the observed times-intersect distributions, where the stiffer materials have a more limited range of times-intersect values. The nonlinear relationship between shear modulus and times-intersect is also consistent with the nonlinear relationship between shear modulus and displacement observed in ARFI imaging~\cite{nightingale_Acoustic_2002}. This stiffness-dependent challenge in distinguishing times-intersect is also observed in each subfigure of Fig.~\ref{fig:tintStat}, the figure that shows $p$-values of hypothesis tests and effect sizes for the times-intersect distributions in numerical simulations involving media of different stiffnesses. Softer materials predominantly show lower $p$-values (i.e. higher statistical significance) and larger effect sizes (i.e. larger differences in means) when compared against other materials, whereas stiffer materials often have higher $p$-values and smaller effect sizes. This suggests that DoPIo times-intersect allows materials with 4 kPa shear modulus differences to be distinguished "more easily" when their shear moduli are lower.
%
% The distribution of time-intersect values (i.e. dynamic range) in Fig.~\ref{fig:tintDistrib} also varied based on the beam sequence used. Both for separate and simultaneous tracking schemes, the overall range of times-intersect across all materials were smaller for F/1.5 ARF pushes compared to F/3.0 pushes. This pattern was observed for times-intersect from displacements tracked using two beams with larger differences in widths (i.e. F/1.5 and F/5.0) versus those with smaller width differences (F/1.5 and F/3.0), as well. Fig.~\ref{fig:tintStat} quantitatively corroborates this observation: times-intersect from F/3.0 pushes (bottom two rows) generally showed lower $p$-values and larger effect sizes across all simulated stiffness pairs compared to F/1.5 pushes (top two rows), as did those using F/1.5 and F/5.0 tFc (even-numbered rows) versus F/1.5 and F/3.0 (odd-numbered rows). This pattern presented both in the separate tracking scheme (left column) as well as for simultaneous tracking (right column). Combined with the observation from the previous paragraph, these results suggest that beam sequences with larger time-intersect distirbutions and tighter time-intersect distributions for simulations of a given stiffness allow for improved differentiation of DoPIo acquisitions from different shear moduli. This property, then, appears to be a function of the beam sequence used - and beam sequences involving wider pushes, simultaneous tracking, and larger differences in beam widths appear to be more effective for differentiating shear moduli.
%
% However, Fig.~\ref{fig:tintStat}(g), the subfigure where times-intersect do not appear to be significantly different for comparisons among stiffer materials, presents an important exception to the preference in beam sequences noted in the previous paragraph. This exception can be attributed primarily to the lack of changes in time-intersect captured using the corresponding beam sequence (i.e. the light purple box plots in Fig.~\ref{fig:tintDistrib}(b)) from materials with shear moduli above 23 kPa. The idea of a PRI-based limit to DoPIo's performance from the previous paragraph may offer one explanation here, as a smaller PRI (i.e. higher PRF) could allow times-intersect to be captured more precisely in materials that quickly recover from shear displacements. However, a higher PRF presents two issues: (1) a higher PRF would also limit the maximum depth of post-ARF imaging, which could be a limiting factor in the clinical setting, and (2) the impact of a higher PRF, as opposed to displacement profile interpolation as applied in this work, on the time-intersect distributions is not yet known. We reserve further investigation of the impact of PRF on DoPIo for future work.
%
% Taking a more holistic view of figures~\ref{fig:tintDistrib} and~\ref{fig:tintStat}, we must note that our comparison of push and track beam sequences for the distribution of times-intersect (and the performance of DoPIo more broadly) implies that a relationship exists between the beam sequence used and the ability to distinguish between different elasticities. This relationship is not trivial since, unlike typical ARFI applications, five F-numbers must be considered: one for the ARF push beam, two when tracking pulses are transmitted, and two to beamform received signals. The results of this work suggest that the choice of beam sequence can have a significant impact on the ability to distinguish between different elasticities, and that this impact can vary depending on the specific beam sequence used. For example, while an F/3.0 push appears to improve distinction between time-intersect distributions versus F/1.5 pushes, tFc that included F/3.0 underperformed those that included F/5.0-tracked displacements. Although it is tempting to conclude that a higher F-number improves DoPIo's performance, previous ARFI imaging experiments of phantoms using F/5.0 pushes did not concentrate enough acoustic energy to induce measurable displacements in stiffer materials. Similarly, when a tFc of F/3.0 versus F/5.0 was used, each tracking beam offered less tracking sensitivity compared to tFc involving F/1.5 beams, resulting in broader distributions of times-intersect (i.e. time-intersect values that were less distinguishable across stiffness). Thus, it may be more appropriate to consider the relative differences in beam widths and the resulting time-intersect distributions rather than the absolute F-number of the beams used. This is consistent with the observed times-intersect distributions, where wider beams and larger differences in beam widths were associated with larger ranges of times-intersect values.
%
% As for the timing of track beams, the results of this work suggest that simultaneous tracking may be more effective for distinguishing between different elasticities than separate tracking. Here, we should recall from Table~\ref{tbl:fnum} that simultaneous tracking was implemented \emph{in silico} by detecting displaced scatterers in Field-II for both beams in a tFc using the same distribution of scatterers. However, a wider transmit F-number was used so that scatterer reflections could be detected by both track beams in and near the focus (e.g. for a tFc of F/1.5 versus F/5.0, track pulses were still transmitted as F/5.0). On the other hand, for simultaneous tracking, different distributions of scatterers were used for the two F-numbers, but the transmit and receive F-numbers for each track beam were matched. Therefore, Fig.~\ref{fig:tintStat} shows that, even though F-numbers between track transmit and receive beams are inconsistent for the simultaneous tracking simulations, the degraded focusing of received backscatter in the simultaneous tracking scheme may be compensated by how both track pulses encode the exact same distribution of scatterers and underlying electronic noise in RF signals. Although the improved performance of simultaneous versus separate tracking was not universal across all imaging parameters, it nevertheless suggests that the act of distinguishing materials of different elasticities using times-intersect benefits from the consideration of the tracking of \emph{truly identical} scatterer motion following an ARF excitation.

\subsection{\emph{In Silico} Modulus Estimation}

As summarized in Table~\ref{tbl:modelr2}, both models based on Eq.~\ref{eq:modelinvsq} and Eq.~\ref{eq:modellog} provided empirical mappings from $t_{\textnormal{int}}$ to shear modulus. The latter form consistently yielded higher coefficients of determination ($R^2 > 0.9$ across most beam sequences) compared to the inverse-square model ($R^2 \approx 0.7 - 0.8$). This is notable since the inverse-square model is most reminiscent of how the time duration of shear wave propagation is inverted and squared to estimate shear elastic modulus in SWEI~\cite{sarvazyan_Shear_1998,palmeri_Ultrasonic_2006}. The improved performance of the model from Eq.~\ref{eq:modellogactual}, then, suggests that the relationship between $t_{\textnormal{int}}$ and shear modulus is more complex. This is reasonable since DoPIo is typically performed on-axis to ARF excitations, where shear waves are not yet fully formed. Thus, $t_{\textnormal{int}}$ is likely informed by a combination of tissue's initial ARF displacement response and subsequent propagation of shearing motion.

In Fig.~\ref{fig:Gdistrib} and Table~\ref{tbl:medianG}, modulus estimates tended to be close to ground truth across all simulated conditions, with median errors well under the step size of 4~kPa. These patterns were exhibited across all beam sequences, with those that used simultaneous tracking, F/3.0 pushes, and F/1.5 versus F/5.0 tracking pairs generally yielding the most accurate estimates. However, the range of modulus estimates for each materials are much larger, leading to much larger interquartile ranges (IQR) in Fig.~\ref{fig:Gdistrib} and median absolute deviations in Table~\ref{tbl:medianG}. Notably, the modulus estimate IQRs from all beam sequences are larger in stiffer materials than in softer ones. This is expected as the precision of modulus estimates is impacted by how, per Eq.~\ref{eq:modelinvsq} and~\ref{eq:modellog}, the same magnitude of change in $t_{\textnormal{int}}$ impacts modulus estimates more for stiffer materials (i.e. faster times-intersect) than for softer ones. Since this heteroskedasticity in modulus estimates directly arises from the process of detecting times-intersect and applying model in the form of Eq.~\ref{eq:modelgeneric}, it is ultimately a fundamental limitation of DoPIo. While this issue may be mitigated to an extent by optimizing times-intersect detection (e.g. using new beam sequences, higher PRF), a nonlinear model - perhaps one that even takes into account a machine learning-based approach where paired displacement profiles are directly inputted instead of $t_{\mathrm{int}}$ - may be needed to improve the precision of modulus estimates in stiffer materials.

% OMITTED FROM SEPTEMBER 2025 DRAFT:
%
% While Fig.~\ref{fig:depth} showed modest bias in $t_{\mathrm{int}}$ over depth, the empirically derived model compensated effectively in generating stiffness estimates. In Table~\ref{tbl:rmseDepth}, as in Fig.~\ref{fig:Gdistrib} the logarithmic model, F/3.0 push, larger difference in track beam widths, and simultaneous tracking all contributed to lower, more consistent RMSE across depths compared to counterpart configurations in each category.

% OMITTED FROM SPRING 2025 DRAFT:
%
% Time-intersect detection, however, is only one "half" of DoPIo: the other half, as shown in Fig.~\ref{fig:Gfit}, is the \emph{relation} of times-intersect to shear elasticity. DoPIo's ability to do so \emph{in silico} is depicted in Fig.~\ref{fig:Gdistrib}, where the distributions of shear moduli estimates are shown for each beam sequence. The median of shear modulus estimates roughly corresponds with the tground truth for all beam sequences. Crucially, the trend between true and estimated shear moduli are heteroskedastic in all beam sequences; stiffer materials are associated with larger variances in estimated shear moduli. This is consistent with the observed times-intersect distributions, where stiffer materials had smaller ranges of times-intersect values that were closer together. Since equations~\ref{eq:modelinvsq} and~\ref{eq:modellog} both relate time-intersect inversely to shear modulus, the same magnitude of change in a time-intersect value is magnified for smaller values (i.e. stiffer materials) as opposed to large ones (i.e. softer materials).
%
% The observed patterns for distinguishing time-intersect in softer versus stiffer materials largely also held true in analyses of $p$-values and effect sizes for Fig.~\ref{fig:Gdistrib}, though the figures were omitted for brevity. Softer materials were generally more likely to be distinguishable from stiffer ones. However, perhaps owing to the wider, more heteroskedastic distributions of modulus estimates, the effect sizes of most comparisons were significantly reduced. This pattern was also observed in the beam sequences with larger differences in beam widths (i.e. F/1.5 and F/5.0) versus those with smaller differences (F/1.5 and F/3.0). It should be noted, however, that the conversion of times-intersect into modulus predictions uses a linear model but is a nonlinear transformation; the translation of distributions and effect sizes, thus, are also nonlinear. This is another important consideration for future work, particularly with the aforementioned considerations for changes in PRI.
%
% We can also compare model fits for all beam sequences and all model types; this was shown in Table~\ref{tbl:modelr2}. Similar to our observations for times-intersect and modulus predictions, $R^2$ values were generally higher for models using an F/3.0 push, tFc with larger differences in track beam widths, and simultaneous tracking. Furthermore, the logarithmic model consistently outperformed the inverse square model. This suggests that the inverse square model, which is more of an analogue to SWEI, may not be the most appropriate model for DoPIo compared to the logarithmic model. The improved appropriateness of the logarithmic model suggests that the relationship between time-intersect and shear modulus is not purely governed by the shear wave speed of the underlying material, but a more complex relationship may exist that links the two quantities.
%
% Next, the \emph{in silico} models created for each depth were evaluated. The results of these models, as well as the corresponding times-intersect distributions, are exemplified in Fig.~\ref{fig:depth} and generally described in Table~\ref{tbl:rmseDepth}. Generally, while times-intersect vary over depth as well as in different materials, Fig.~\ref{fig:depth}(b) shows that these differences may be counteracted by developing different models (i.e. different coefficients for equations~\ref{eq:modelinvsq} and~\ref{eq:modellog}) for different depths. Furthermore, the RMSE of modulus predictions in Table~\ref{tbl:rmseDepth} were generally more consistent over depth (i.e. more consistent across columns for a given row) when wider pushes, simultaneous beamforming of displacements-tracking pulse-echoes, tFc with larger width separations, and/or the model using logarithmic transform were used. It should be noted that the exact relationships between these four parameters, as well as changes therein over depth, were not consistent across all permutations. Still, all beam sequence and model variants led to modulus estimate RMSEs of less than 4.0 kPa, the step size of simulated material elasticities. This suggests that the DoPIo technique is capable of estimating shear moduli with a precision of at least 4.0 kPa. 
%
% In summary, \emph{in silico} acquisitions of DoPIo based on FEM simulations for materials of known elasticities resulted in times-intersect that could be distinguished based on the stiffness of the ARF-interrogated material. Using equations~\ref{eq:modelinvsq} and~\ref{eq:modellog}, a linear regression model was developed to relate times-intersect to modulus for a given depth. This model, then, could be evaluated on times-intersect from additional displacement profile pairs. The results of this evaluation showed that the model was able to predict shear moduli with a precision of 4.0 kPa, and that the model fit was generally improved by using an F/3.0 push, simultaneous tracking, and tFc with larger differences in beam widths. The logarithmic transform-based model based on Eq.~\ref{eq:modellog} also generally outperformed the inverse square model (i.e. Eq.~\ref{eq:modelinvsq}) for all beam sequences. This relationship was evaluated both at 25.00 mm, the focal depth of the ARF push and track beams, as well as for a broader range of depths where displacements were tracked.

\subsection{Experimental Acquisitions}

Fig.~\ref{fig:cirsImage} shows that DoPIo reproduces the expected qualitative contrast in a calibrated phantom; inclusions that are nominally stiffer appear with smaller $t_{\textnormal{int}}$ and higher DoPIo-estimated modulus relative to the background, and vice versa. This suggests that, at least qualitatively, DoPIo as implemented using FEM-derived empirical models can use time-intersect to estimate shear moduli. Furthermore, although the model was developed solely in homogeneous media, the results of this work suggest that the model can be used to estimate shear moduli in heterogeneous media as well. DoPIo was also able to delineate inclusions at depths away from the ARF focus, though modulus estimates were less accurate there. 

Fig.~\ref{fig:cirsBoxplot}(a)-(f) also show that the morphology of DoPIo images appear largely consistent across different amplitudes of ARF excitations. In Fig.~\ref{fig:cirsBoxplot}(g), DoPIo modulus estimates preserved the correct rank ordering of features with different stiffnesses. Notably, modulus estimates were consistent across all ARF amplitudes in the Type 1 inclusion, and for all but one pair of comparisons across push powers in the Type 2 inclusion. However, modulus estimates in the Type 3 and 4 inclusions were significantly different for most comparisons of push powers with differences of 10 V or greater. For the background material, whose elasticity falls in between those for all four inclusions, modulus predictions were not statistically significant across push powers with a difference under 10 V - even though these comparisons had a fivefold larger sample size than the comparison of inclusions. For all materials except the Type 1 inclusion, an increase in push power coincided with increases in modulus estimates up to some asymptotic value.

Therefore, DoPIo may offer some degree of robustness to variations in ARF excitation amplitude. This is notable since ARF push amplitudes, which are generally unknown, are typically considered to be confounding factors for quantitative elastography techniques that require correction~\cite{deng_Analyzing_2015, chen_ShearWave_2023}. However, DoPIo appears to be robust to this factor to a certain extent. Since the displacement profiles used to compute $t_{\textnormal{int}}$ are normalized by the maximum of that of the narrower beam, the SNR of displacement profiles - and, therefore, the reliability of $t_{\textnormal{int}}$ detection - are degraded where ARF displacement magnitudes are small to the point where displacement signals are overwhelmed by jitter~\cite{mcaleavey_Estimates_2003,dhanaliwala_Assessing_2012}. This issue can take place (1) in stiffer materials, (2) following low-amplitude ARF pushes, (3) at depths away from the ARF push focus, and (4) where the wider tracking beam has a much lower sensitivity to displacements than the narrower one (i.e. when the difference in F-numbers is large). Since DoPIo always uses the first intersection following both peak displacements, this means that displacement profile pairs that are noisy due to any of these four conditions lead to times-intersect occurring later than they "should", leading to modulus understimation. Thus, the performance of DoPIo-based modulus estimation is independent of ARF push power as long as the displacement profiles have sufficient SNR to detect $t_{\textnormal{int}}$ accurately.

Fig.~\ref{fig:liver}(a)-(c) demonstrate that DoPIo can also depict elastic heterogeneities in tissue, though Fig.~\ref{fig:liver}(d) shows systematic underestimation relative to SWEI-based measurements from the clinical scanner. Note that the SWEI-based measurements are solely used in this study as a comparison metrics in the absence of ground truth values. Still, previous work shows that SWEI-based shear wave velocity measurements of the same medium can vary by around 10\% based on scanner systems, transducers, imaging and outlier-rejection settings, and other implementation-specific parameters~\cite{palmeri_Radiological_2021} - a margin of error that the DoPIo results exceed.  This also suggests that future work is needed to refine the empirical model to improve the accuracy of non-simulated DoPIo-based modulus estimates. Thus, DoPIo was able to generate quantitative images of tissue stiffness estimates in both \emph{in vitro} and in biological tissue, suggesting its potential for clinical applications.

% OMITTED FROM SEPTEMBER 2025 DRAFT - replace with last sentence of previous paragraph
%
%Lastly, Fig.~\ref{fig:kidney} shows that DoPIo can also differentiate tissue structures \emph{in vivo}. Thus, DoPIo was able to generate quantitative images of tissue stiffness estimates in both \emph{ex vivo} and \emph{in vivo} settings, suggesting its potential for clinical applications.

% OMITTED FROM SPRING 2025 DRAFT:
%
% \ subsection{\emph{In Vitro} Acquisitions}
%
% Representative examples of \emph{in vitro} acquisitions of DoPIo are shown in Fig.~\ref{fig:cirsImage}, whereas a more quantitative analysis of modulus estimates are shown in Fig.~\ref{fig:cirsBoxplot}. In the former figure, times-intersect distributions in the CIRS phantom show that the times-intersect values are visually distinct between the inclusion and background material, and that the regions where those differences present are delineated by the outline of the inclusion. Times-intersect near the focal depth of 25.00 mm are also noticeably different across the three inclusions of increasing stiffness. This trend is consistent with the simulation results, as the softer Type 1 inclusion (i.e. 2.2 kPa) showed slower times-intersect (i.e. higher values, brighter colors) than the stiffer Type 3 inclusion (i.e. 16.3 kPa), with the background and Type 2 inclusions falling in between those extremes. The images of shear modulus estimates also show a similar trend, with the Type 1 inclusion showing lower estimated shear moduli than the background material or stiffer inclusions. This suggests that, at least qualitatively, DoPIo as implemented using FEM-derived empirical models can use time-intersect to estimate shear moduli. Furthermore, although the model was developed solely in homogeneous media, the results of this work suggest that the model can be used to estimate shear moduli in heterogeneous media as well. This is an important consideration for the clinical utility of DoPIo, as tissue is highly heterogeneous.
%
% Interestingly, Fig.~\ref{fig:cirsBoxplot}(a)-(f) also show that the morphology of DoPIo images appear largely consistent across different amplitudes of ARF excitations. While the modulus estimate of the background appear to get slowly brighter with increasing push powers (especially between 45 V and 50 V), all images of modulus estimates show elastic differences in and out of the inclusion near the focus. This pattern is also visible in Fig.~\ref{fig:cirsBoxplot}(g), as the distributions of modulus estimates for stiffer inclusions appear to slightly increase with an increase in push power amplitude. For example, DoPIo modulus values for the Type 3 inclusion underestimated the true value when lower push amplitudes were used, but those using 55 V and 60 V pushes resulted in interquartile ranges that encompassed the ground truth. More generally, differences in modulus estimates were only statistically significant in stiffer inclusions between push powers differences of 15 V or greater, whereas the softest Type 1 inclusion did not have significant differences across any of the evaluated push powers. For the background material, whose manufacturer-reported elasticity falls in between those for all four inclusions, modulus predictions were not statistically significant across push powers with a difference under 10 V - even though these comparisons had a fivefold larger sample size than the comparison of inclusions. This suggests that DoPIo may offer some degree of robustness to variations in ARF excitation amplitude. This is an important consideration for the clinical utility of DoPIo, as variations in ARF excitation amplitude can occur due to differences in system settings or patient anatomy.
%
% Still, shear modulus values estimated in the calibrated phantom did not match manufacturer-reported values. The image-sampling regions in Fig.~\ref{fig:cirsBoxplot}(g) were located near the focus of ARF push and track beams, essentially acting as best-case scenarios for delineating elastic differences. DoPIo especially underestimated the elasticity of Type 3 and 4 inclusions even in those idealized depths, with the estimated stiffness of the Type 4 inclusions at maximum push amplitude being half of the expected value. Figures~\ref{fig:cirsImage} and~\ref{fig:cirsBoxplot}(a)-(f) show that this error is exacerbated away from the focus (i.e. around 15-20 mm and 30-35 mm), where elasticity estimates appear to severely underestimate the true stiffness of the background material. It should be noted, however, that times-intersect from the experimentally acquired images are also much higher than times-intersect for similar stiffnesses and depths as seen in Fig.~\ref{fig:depth}. This suggests that the inaccuracy of experimental estimations of elasticity using DoPIo are not solely due to inaccuracies in the model; instead, the distribution of times-intersect \emph{in silico} is inconsistent with those \emph{in vitro}. This point and its implication is further discussed in the final subpart of this discussion.
%
% \ subsection{\emph{Ex Vivo} and \emph{In Vivo} Acquisitions}
%
% Similarly to Fig.~\ref{fig:cirsImage}, Fig.~\ref{fig:liverImage} shows an image of DoPIo-estimated shear elasticity in \emph{ex vivo} biological tissue. The image clearly shows a difference in elasticity between the stiffer gelatin inclusion, outlined in white based on B-mode, versus a porcine liver background. However, Fig.~\ref{fig:liverBoxplot} shows that DoPIo-derived elasticity estimates are significantly lower than those from a clinically-used implementation of SWEI. Similar to Fig.~\ref{fig:cirsBoxplot}(g), the estimated shear modulus of the liver background material also slightly increases as a function of push power, though the differences were not statistically significant.
%
% We must recall that the SWEI-based measurement of liver and inclusion stiffness is solely used in this study as a comparison metric in the absence of ground truth values. Notably, the clinical scanner used a different scanner system, transducer, and imaging parameters, and the analysis of image quality is complicated by how the Acuson Sequoia and other clinical scanners apply black-box filters to smooth image outputs and hide low-quality data. Due to these and other factors, SWEI-based shear wave velocity measurements are expected to vary by around 10\% based on system manufacturer and other implementation-specific parameters~\cite{palmeri_Radiological_2021}. While these issues may lead to some of the differences in modulus estimates between DoPIo and SWEI, the roughly twofold difference in DoPIo- and SWEI-based stiffness estimates suggest that there are other factors at play - perhaps those that relate to the same issue of modulus inaccuracy described in the previous subsection. Thus, we should consider DoPIo to be a promising technique to generate qualititative metrics for tissue stiffness, but one that requires further refinement before its elasticity estimates can be used for clinical decision-making.
%
% Finally, Fig.~\ref{fig:kidney} shows an example image of DoPIo-estimated shear elasticity, \emph{in vivo}, in a porcine kidney. The image shows a clear difference in elasticity between the renal cortex and medulla, with the cortex appearing softer than the medulla. It should be noted that the true stiffness of various parts of the renal parenchyma are not well-established in the literature, in part due to a lack of standardization in how to account for the elastic anisotropy of the organ. Still, the ability to quantitatively differentiate structural differences of a heterogeneous organ \emph{in vivo} is a promising step towards the clinical utility of DoPIo. Future work is needed to validate the accuracy of DoPIo's modulus estimates in \emph{in vivo} settings, to explore the potential for using DoPIo to assess the mechanical properties of other organs and tissues, and to assess the impact of mechanical anisotropy on DoPIo's performance.

\subsection{Future Directions}

A central limitation is the mismatch between simulated and experimental time-intersect distributions. While \emph{in silico} studies suggested close agreement between estimated and true moduli, phantom and tissue results showed systematic underestimation. This discrepancy stems from known discrepancies between the shapes of simulated and experimental on-axis displacement profiles, the cause of which is still under investigation. Future efforts should focus on closing the gap, either through the creation of a larger dataset of experimental acquisitions of custom-made phantoms or through data-driven strategies for domain adaptation. Once these issues are addressed, additional topics such as the evaluation of DoPIo's performance in viscoelastic media or additional beamforming techniques may be investigated.

% previous sentence omitted from SEPTEMBER 2025 DRAFT:
%
% This discrepancy may stem from idealized assumptions for how ARF momentum transfer and displacement responses are modeled in LS-DYNA as well as how how that is represented in Field-II.

Likewise, quantitative calibration also remains incomplete. Although DoPIo consistently showed the correct relationship between softer and stiffer media, it underestimated moduli for stiffer materials. Empirical models based on logarithmic fitting reduced error but did not eliminate this bias. Future efforts should refine the mapping from time-intersect to modulus, either through improved physical modeling or data-driven regression. Other technical challenges include optimizing beam sequences and addressing depth-dependent bias. These are secondary to the issues of accuracy in stiffer media and precision, but they remain relevant for clinical translation. Machine learning and domain adaptation approaches may further enhance robustness by directly linking simulated and experimental distributions.

In summary, DoPIo has demonstrated feasibility for qualitative stiffness imaging across simulations, phantoms, and tissue. The next priorities are: (i) refining models to close the gap between simulation and experiment, (ii) establishing quantitative calibration across diverse tissues and pathologies, and (iii) exploring data-driven approaches to improve the accuracy and precision of modulus estimates. Addressing these challenges will be critical to translating DoPIo into a clinically reliable imaging modality.

% OMITTED FROM SPRING 2025 DRAFT:
%
% A critical limitation of DoPIo is that its ability to quantify and distinguish materials with different shear elasticities depends on the quality of the empirically derived model. This, in turn, depends on whether the FEM-derived model and the pipeline of \emph{in silico} model development are reflective of experimental observations. Therefore, it is critical that simulated times-intersect vary over stiffness and depth in a manner that resembles those from experimental acquisitions. However, there is a clear discrepancy in DoPIo's performance \emph{in silico} versus in physical experiments. Unlike in Fig.~\ref{fig:depth}, Fig.~\ref{fig:cirsImage} and~\ref{fig:cirsBoxplot} show that both times-intersect and modulus estimates substantially varied in the background material of the tissue-mimicking phantom within versus past 5 mm from the focal depth of 25 mm. While times-intersect are clearly different between the inside and outside of inclusion borders, the resulting stiffness estimates are still incorrect. This inaccuracy suggests two broad areas where DoPIo may be further developed: (1) the general process of inducing displacements and tracking displacements to detect times-intersect that are more amenable to DoPIo-based modulus prediction, as well as (2) the specific workflow used to convert paired, simulation-based, ultrasonically-tracked displacement estimates into empirical models that enable shear modulus estimation.
%
% \ subsubsection{On Optimizing Times-Intersect Distributions}
%
% The first area of improvement concerns the process of inducing displacements and tracking displacements to counteract DoPIo's degraded ability to accurately distinguish materials with stiffer elasticities \emph{in silico} and \emph{in vitro}. The results of this work suggest that, generally, an F/3.0 push beam followed by simultaneous tracking using F/1.5 and F/5.0 beams best distinguishes displacement acquisitions from media with different elasticities. However, the superiority of that beam sequence did not hold true when times-intersect distributions were compared between some stiffer materials in Fig.~\ref{fig:tintStat}, as well as in some cases of modulus estimates for stiffer materials. We can recall that an F/3.0 push beam generally produced larger ranges of times-intersect values compared to F/1.5 pushes and tFc with larger differences in beam widths were associated with improved separations of times-intersect by stiffness. Given this, future work could investigate if DoPIo's performance could be improved using other beamforming approaches for ARF pushes or by using tracking aperture pairs whose differences in beam widths can be further exaggerated such that times-intersect are more distinct across different elasticities. 
%
% Similarly, the aforementioned push and track beam configuration deviated from its general position as the "best" one in the RMSE of modulus estimates above and below the focus in Table~\ref{tbl:rmseDepth}. This also suggests that the \emph{in silico} dataset, alone, presents a need to improve how times-intersect are detected and related by elasticity in depths away from the push and track focal depth. A preliminary investigation showed that, using ARF excitations with a Bessel function-based apodization~\cite{feng_Shear_2021}, DoPIo times-intersect and modulus estimates were more consistent over depth compared to focused beam-based implementations \emph{in silico}~\cite{yokoyama_Bessel_2023}. Future work may also investigate the use of Bessel function-based beams in physical experiments to see whether this, or other beamforming techniques, can improve the consistency of times-intersect and modulus estimates over depth.
%
% \ subsubsection{On Improving the DoPIo Model Process}
%
% The other area of improvement concerns the process of converting paired, simulation-based, ultrasonically-tracked displacement estimates into empirical models that enable shear modulus estimation. First, we can examine how Field-II is used to simulate the ultrasonic tracking of a large number of displaced scatterers, thereby creating an \emph{in silico} example dataset of times-intersect. Given the Rayleigh model of sub-wavelength scatterers, this requires that a minimum number of scatterers exist within a resolution cell volume for fully-developed speckles to realize in ultrasound backscatter signals \cite{wagner_Statistics_1983,wagner_Fundamental_1988,trahey_Properties_1988}. Since the signal-to-noise ratio of envelope-detected images plateaued at roughly 11 per resolution cell for B-mode imaging, this scatterer density is typically considered a reasonable default value for determining the density of scatterers in a Field-II simulation. However, DoPIo relies on the detection of differences in scatterer motion between two resolution cells. We identified 25 scatterers per resolution cell as a reasonable balance where times-intersect were detectable but total simulation times were reasonable on the supercomputer cluster provided by the authors' institution. However, this analysis was limited in scale, and a more systematic analysis of acoustic wave simulators for DoPIo may also be warranted. 
%
% Besides issues that are fundamental to computationally simulated ultrasonic tracking, we can consider how the FEM-derived DoPIo models for stiffness predictions get the general trend correct, but additional data-driven approaches are required to improve the accuracy of modulus estimates. This effectively makes the problem of improving DoPIo into a domain adaptation problem, where the goal is to adapt the model to a new domain (i.e. experimental data) that is different from the training domain (i.e. \emph{in silico} data). This could be accomplished by using a small amount of labeled data from the new domain to fine-tune the model, or by using unsupervised domain adaptation techniques to align the distributions of the two domains.
%
% For example, Richardson~\emph{et al.}~\cite{richardson_Quantitative_2024} showed that a fully-connected neural network (FCNN) trained on displacement profiles can be relevant for this task. Like DoPIo, this quantitative viscoelastic response (QVisR) elastography technique estimates shear elasticity using the same simulation pipeline for generating \emph{in silico} examples of on-axis-tracked displacements. While QVisR requires two ARF excitations and relies on training data whose FEM material property models were more advanced than those for DoPIo, it also produced an empirical model that relates displacement profile-based information to material property estimates - one that, like DoPIo, also struggled to accurately estimate the stiffness of elastic heterogeneities over multiple depths in physical experiments. However, QVisR's accuracy was improved using a two-step training process where initial FCNN training using \emph{in silico} displacement profiles followed by transfer learning using \emph{in vitro} data from calibrated phantoms. While preliminary work has demonstrated that DoPIo using support vector regression in lieu of linear regression is feasible~\cite{rahmouni_Machine_2020}, the impact of DoPIo's performance using more complex models and transfer learning remains to be investigated.
%
% To simplify this investigation, one may find it tempting to develop an analytical solution that directly derives elasticity based on simplified representations of a push beam, paired track beams, and times-intersect. However, existing solutions of the elastic wave equation tailored to ARF-based applications, such as those by Sarvazyan \emph{et al.}~\cite{ssarvazyan_Shear_1998,sarvazyan_Acoustic_2021}, involve assumptions that do not obviously allow for its extension to DoPIo, and the authors of this work are not aware of analytical solutions or pertinent Green's functions that could be used to derive an analytical expression that relates times-intersect to shear modulus as a replacement of empirically derived models.

